\documentclass[11pt,letterpaper]{article}
\usepackage[margin=0.82in,headheight=14pt]{geometry}
\usepackage[T1]{fontenc}
\usepackage{lmodern,amsmath,amssymb,amsthm,microtype}
\usepackage{xcolor,hyperref,fancyhdr,listings,enumitem,booktabs,tabularx}
\definecolor{ink}{HTML}{18364A}
\definecolor{wash}{HTML}{F3F6F8}
\hypersetup{colorlinks=true,linkcolor=ink,urlcolor=ink,pdftitle={Interior smoothness of weak heat solutions},pdfauthor={Fluid PDEs in Lean Foundation}}
\pagestyle{fancy}\fancyhf{}
\fancyhead[L]{\small\textsc{Fluid PDEs in Lean Foundation}}
\fancyhead[R]{\small Interior regularity}
\fancyfoot[C]{\small\thepage}
\setlength{\parindent}{0pt}\setlength{\parskip}{6pt}
\setlist{leftmargin=*,itemsep=5pt,topsep=4pt}
\newtheorem{theorem}{Theorem}
\newcommand{\R}{\mathbb R}
\newcommand{\D}{\mathcal D}
\newcommand{\Lloc}{L^1_{\mathrm{loc}}}
\lstset{basicstyle=\ttfamily\footnotesize,backgroundcolor=\color{wash},
  frame=single,rulecolor=\color{black!15},framesep=7pt,
  columns=fullflexible,keepspaces=true,breaklines=true,showstringspaces=false,
  escapeinside={(*@}{@*)}}
\setlength{\emergencystretch}{2em}
\begin{document}
{\color{ink}\LARGE\bfseries Interior smoothness of weak heat solutions\par}
{\large English theorem, checked Lean statement, and proof outline\par}
{\small Prepared 18 October 2026\quad\textbullet\quad Homogeneous equation, diffusivity one}

\section*{1. Precise mathematical statement}
Write a space--time point as $z=(t,x)\in\R\times\R^n$, and let
$P=\partial_t-\Delta_x$. All smoothness below is joint in time and space.

\begin{theorem}[Interior regularity on an open cylinder]
Let $n\geq1$, let $\Omega\subseteq\R^n$ be open, let $T>0$, and set
$Q=(0,T)\times\Omega$. Suppose $u\in\Lloc(Q)$ and
\begin{equation}\label{eq:weak}
  \int_Q u(t,x)\bigl(\partial_t\varphi(t,x)+\Delta_x\varphi(t,x)\bigr)
  \,dt\,dx=0
  \qquad\text{for every }\varphi\in C_c^\infty(Q).
\end{equation}
Then there is a unique $v\in C^\infty(Q)$ such that $u=v$ almost everywhere on
$Q$. Moreover,
\[
  \partial_t v(t,x)=\Delta_x v(t,x)\qquad\text{at every }(t,x)\in Q.
\]
Uniqueness is among smooth representatives of this same almost-everywhere
class on $Q$.
\end{theorem}

Equation~\eqref{eq:weak} is precisely $Pu=0$ in distributions. Its plus sign
is intentional: integration by parts gives
$\langle Pu,\varphi\rangle=-\int u(\partial_t\varphi+\Delta_x\varphi)$.
The formalization assumes local integrability, without pointwise differentiability
of $u$.

\textbf{Stronger proved form.}
The same conclusion holds for any open $Q\subseteq\R^{n+1}$; a product
structure is unnecessary. The cylinder theorem follows from this result.

\section*{2. Domain assumptions and boundary scope}
Only openness of $\Omega$ is needed. It need not be bounded, connected, or
have a Lipschitz boundary. Each point of $Q$ has a neighborhood whose closure
lies in $Q$, and the proof works in that neighborhood.

The theorem gives no extension of smoothness to $\partial\Omega$, $t=0$, or
$t=T$, and no uniform derivative estimates as these sets are approached.
An irregular boundary may affect such boundary conclusions without affecting
interior smoothness. A graph such as $y=\sqrt{x}$ has an unbounded slope in
that coordinate system; this alone does not settle whether the surrounding
domain is Lipschitz, since boundary charts can be rotated.

The representative qualification is essential: altering $u$ on a null set
can destroy pointwise continuity while preserving~\eqref{eq:weak}.

\newpage
\section*{3. Proof outline matching the formalization}
The following explains the checked proof; the complete formal proof is in the
repository modules listed on page~\pageref{sec:verification}.

\textbf{Step 1: Localize and regularize.}
Fix $z_0\in Q$. Choose a compact neighborhood $K\subset Q$ and a smooth
cutoff $\chi$ supported in $\operatorname{int}K$, with $\chi=1$ near $z_0$.
Extend $u\mathbf1_K$ by zero and convolve with normalized smooth bumps whose
radii tend to zero. The resulting $u_j$ are smooth, converge to the localized
input in $L^1$ and almost everywhere, and are eventually caloric across the
whole compact set $\operatorname{supp}\chi$. The last assertion follows by
transferring derivatives to test kernels in the distributional identity.

\textbf{Step 2: Prove the fundamental-solution identity.}
Construct the causal Gaussian
\[
 H(t,x)=\begin{cases}
 (4\pi t)^{-n/2}\exp\!\bigl(-|x|^2/(4t)\bigr),&t>0,\\
 0,&t\leq0.
 \end{cases}
\]
It is locally integrable, smooth away from $(0,0)$, has unit spatial mass
at positive times, and approaches the spatial Dirac mass initially. The proved
identity, for smooth compactly supported $\varphi$, is
\[
  \int_{\R^{n+1}} H(z)(\partial_t\varphi+\Delta_x\varphi)(z)\,dz
  =-\varphi(0).
\]
Equivalently, $PH=\delta_0$. The proof uses explicit Gaussian derivatives,
integration by parts, the initial limit, and Fubini.

\textbf{Step 3: Remove derivatives of the solution.}
Apply the fundamental solution to $\chi u_j$. The cutoff product rule and
spatial integration by parts give, for $z$ near $z_0$,
\begin{align*}
 u_j(z)={}&\int H(z-y)u_j(y)
                   (\partial_t\chi+\Delta_x\chi)(y)\,dy\\
 &-2\sum_{i=1}^n\int(\partial_{x_i}H)(z-y)u_j(y)
                                      \partial_{x_i}\chi(y)\,dy.
\end{align*}
All derivatives on the right act on the cutoff or kernel.

\textbf{Step 4: Pass to the weak solution.}
The supports of the differentiated cutoffs avoid $z_0$. Shrink its neighborhood
so that each kernel can be replaced, on all relevant differences $z-y$, by a
fixed compactly supported smooth kernel. The right side then defines a smooth
function $v$, since it is a sum of convolutions of integrable functions with
smooth compact kernels. Strong $L^1$ convergence passes the representation to
$v$, and almost-everywhere convergence identifies $v$ with $u$ locally.
Thus the smoothness of the limit comes from a fixed convolution representation.

\textbf{Step 5: Glue and recover the equation.}
Local smooth representatives agree on overlaps: continuous functions equal
almost everywhere on an open set agree everywhere there. Glue them, using a
countable subcover to preserve global almost-everywhere equality. The proved
classical/distributional equivalence for smooth functions gives $Pv=0$
pointwise. The same continuity argument gives uniqueness.

\newpage
\section*{4. Exact Lean statement and proof application}
The code below is the example extracted from the repository documentation and
compiled successfully. Here \texttt{Space (n + 1)} is
\texttt{Fin (n + 1)} $\to\R$: coordinate $0$ is time and coordinates
\texttt{i.succ} are spatial. The measure is space--time Lebesgue measure.

\begin{lstlisting}
import PDE.HeatInteriorRegularity

open MeasureTheory Set
open scoped ContDiff BigOperators

namespace PDE.HeatInteriorSpecification

example (n : (*@$\mathbb N$@*)) (hn : 0 < n)
    ((*@$\Omega$@*) : Set (Space n)) (h(*@$\Omega$@*) : IsOpen (*@$\Omega$@*))
    (T : (*@$\mathbb R$@*)) (hT : 0 < T)
    (u : Space (n + 1) (*@$\to$@*) (*@$\mathbb R$@*))
    (hu : IsDistributionalHeatSolution (heatCylinder (*@$\Omega$@*) T) u) :
    (*@$\exists$@*) v : Space (n + 1) (*@$\to$@*) (*@$\mathbb R$@*),
      ContDiffOn (*@$\mathbb R$@*) (*@$\infty$@*) v (heatCylinder (*@$\Omega$@*) T) (*@$\wedge$@*)
      u =(*@${}^{\mathrm m}$@*)[volume.restrict (heatCylinder (*@$\Omega$@*) T)] v (*@$\wedge$@*)
      ((*@$\forall$@*) z (*@$\in$@*) heatCylinder (*@$\Omega$@*) T,
        coordDeriv 0 v z = spaceTimeLaplacian v z) (*@$\wedge$@*)
      ((*@$\forall$@*) w : Space (n + 1) (*@$\to$@*) (*@$\mathbb R$@*),
        ContDiffOn (*@$\mathbb R$@*) (*@$\infty$@*) w (heatCylinder (*@$\Omega$@*) T) (*@$\to$@*)
        u =(*@${}^{\mathrm m}$@*)[volume.restrict (heatCylinder (*@$\Omega$@*) T)] w (*@$\to$@*)
        EqOn w v (heatCylinder (*@$\Omega$@*) T)) :=
  heat_interior_smooth hn h(*@$\Omega$@*) hT hu

end PDE.HeatInteriorSpecification
\end{lstlisting}

\texttt{IsDistributionalHeatSolution Q u} contains exactly local integrability
and the test identity~\eqref{eq:weak}. The output includes joint smoothness,
almost-everywhere equality, the pointwise equation, and uniqueness on the domain.
The function $v$ is defined on all of $\R^{n+1}$ in Lean, but its values outside
$Q$ are unrestricted by this conclusion.

\newpage
\section*{5. Worked examples}
In one spatial dimension, $u(t,x)=x^2+2t$ satisfies
\[
  \partial_tu=2=\partial_x^2u.
\]
The checked polynomial example establishes the distributional equation using
the proved smooth classical/weak equivalence. It also checks the derivative
and sign conventions used throughout the development.

In two spatial dimensions, consider the open cusp
\[
 \Omega_{\mathrm{cusp}}=
 \{(x_1,x_2):0<x_1<1,\ |x_2|<x_1^2\}.
\]
The second checked example proves openness directly from these strict
inequalities, then applies the theorem to any distributional heat solution on
$(0,1)\times\Omega_{\mathrm{cusp}}$. No boundary regularity hypothesis enters
that application. The example does not formalize a separate geometric theorem
about non-Lipschitzness of this cusp.

\section*{6. Formal implementation and verification}\label{sec:verification}
\begin{tabularx}{\linewidth}{@{}p{0.39\linewidth}X@{}}
\toprule
\textbf{Module family} & \textbf{Role}\\\midrule
\texttt{HeatDistribution} & Weak equation, a.e. invariance, classical equivalence.\\
\texttt{HeatMollification},\newline\texttt{HeatApproximation},\newline\texttt{HeatApproximationL1} & Smooth caloric approximations and strong $L^1$ convergence.\\
\texttt{HeatGaussian*} & Gaussian calculus, initial limit, causal smoothness, integrability, and fundamental identity.\\
\texttt{HeatCutoffRepresentation},\newline\texttt{HeatKernelExtension} & Local representation and fixed smooth replacement kernels.\\
\texttt{HeatLocalRegularity},\newline\texttt{HeatInteriorRegularity} & Local smooth representatives, gluing, and the final theorems.\\
\bottomrule
\end{tabularx}

\textbf{Checked baseline.} Commit \texttt{4324148}, prepared 18 October 2026.
The full library and five example modules passed
\texttt{bash scripts/check.sh}. All 560 named theorem roots were audited,
including 82 added heat roots. Only \texttt{propext}, \texttt{Classical.choice},
and \texttt{Quot.sound} occur; there are no \texttt{sorry} proofs or new axioms.
This is a Lean kernel build and transitive axiom audit; an independent kernel
replay was not performed.

The documented example passed
\texttt{python3 scripts/check-heat-interior-statement.py}.
Lean is pinned to \texttt{4.34.0-rc2}; Mathlib is pinned to
\texttt{85e3a25e006c35636f0e53b0e9296caca2685bc0}.
Compilation uses a 3000 MiB Lean budget and a sampled 3500 MiB process-tree
RSS guard. No subagents were used.

\textbf{Solution-class scope.} The theorem uses the scalar distributional
predicate above. An adapter from the older Bochner
\texttt{IsIntegratedWeakHeatSolution} API remains separate future work.

\small
\textbf{Source:} \href{https://github.com/edcrfv/Fluid-PDEs-in-Lean-Foundation}{GitHub repository}.\par
Branch: \path{examples/heat-interior-regularity}. Exact proof:
\path{PDE/HeatInteriorRegularity.lean}; examples:
\path{Checks/HeatInteriorExamples.lean}. Repository licensing and attribution
are recorded in \texttt{LICENSING.md}.
\end{document}
